% !TEX program = lualatex
% Generated by new_lecture_note.sh
\documentclass[11pt]{article}

\usepackage{fontspec}
\usepackage{microtype}
\usepackage{mathtools,amssymb,amsfonts,amsthm,bm}
\usepackage{enumitem}
\usepackage{booktabs,array,xcolor}
\usepackage{geometry,fancyhdr,setspace}
\usepackage[hidelinks]{hyperref}
\usepackage[nameinlink,noabbrev]{cleveref}

\geometry{margin=1in,headheight=15pt}
\setstretch{1.12}
\setlength{\parindent}{0pt}
\setlength{\parskip}{0.55em plus 0.12em minus 0.08em}
\setlength{\abovedisplayskip}{0.8em plus 0.2em minus 0.1em}
\setlength{\belowdisplayskip}{0.8em plus 0.2em minus 0.1em}
\allowdisplaybreaks[3]
\setlist{leftmargin=*,topsep=0.35em,itemsep=0.15em,parsep=0pt}

% ---------- Metadata ----------
\newcommand{\studentname}{Keshav Ramamurthy}
\newcommand{\coursename}{Stat 150}
\newcommand{\lecturenumber}{01}
\newcommand{\lecturetitle}{Lecture 01}
\newcommand{\lecturedate}{\today}

% ---------- Reusable notation ----------
\newcommand{\N}{\mathbb{N}}
\newcommand{\Z}{\mathbb{Z}}
\newcommand{\Q}{\mathbb{Q}}
\newcommand{\R}{\mathbb{R}}
\newcommand{\C}{\mathbb{C}}
\newcommand{\F}{\mathbb{F}}
\newcommand{\K}{\mathbb{K}}
\newcommand{\E}{\mathbb{E}}
\newcommand{\Prob}{\mathbb{P}}
\newcommand{\1}{\mathbf{1}}
\newcommand{\eps}{\varepsilon}
\newcommand{\dd}{\,\mathrm{d}}
\newcommand{\abs}[1]{\left\lvert#1\right\rvert}
\newcommand{\norm}[1]{\left\lVert#1\right\rVert}
\newcommand{\ip}[2]{\left\langle#1,#2\right\rangle}
\newcommand{\set}[1]{\left\{#1\right\}}
\newcommand{\given}{\,\middle\vert\,}
\newcommand{\indep}{\perp\!\!\!\perp}
\newcommand{\lto}{\longrightarrow}
\newcommand{\deriv}[2]{\frac{\mathrm{d}#1}{\mathrm{d}#2}}
\newcommand{\pderiv}[2]{\frac{\partial#1}{\partial#2}}
\DeclareMathOperator{\Aut}{Aut}
\DeclareMathOperator{\Hom}{Hom}
\DeclareMathOperator{\Ker}{ker}
\DeclareMathOperator{\im}{im}
\DeclareMathOperator{\rank}{rank}
\DeclareMathOperator{\nullity}{nullity}
\DeclareMathOperator{\Span}{span}
\DeclareMathOperator{\Var}{Var}
\DeclareMathOperator{\Cov}{Cov}

% ---------- Note environments ----------
\newtheoremstyle{notestyle}{0.7em}{0.7em}{\itshape}{}{}{\bfseries}{.5em}{}
\theoremstyle{notestyle}
\newtheorem{theorem}{Theorem}[section]
\newtheorem{lemma}[theorem]{Lemma}
\newtheorem{proposition}[theorem]{Proposition}
\newtheorem{corollary}[theorem]{Corollary}
\theoremstyle{definition}
\newtheorem{definition}[theorem]{Definition}
\newtheorem{example}[theorem]{Example}
\newtheorem{remark}[theorem]{Remark}
\newenvironment{claim}{\par\noindent\textbf{Claim.}\ }{\par}
\newenvironment{idea}{\par\noindent\textit{Idea.}\ }{\par}
\newenvironment{proofsketch}{\par\noindent\textbf{Proof sketch.}\ }{\par}

% ---------- Header ----------
\pagestyle{fancy}
\fancyhf{}
\fancyhead[L]{\coursename}
\fancyhead[C]{\lecturetitle}
\fancyhead[R]{\studentname}
\fancyfoot[C]{\thepage}

\title{\vspace{-1.5em}\textbf{\coursename -- \lecturetitle}}
\author{\studentname}
\date{\lecturedate}

\begin{document}
\maketitle

\section{Markov Chain}
Markov Chain: Current state is enough information 
$$\mathbb{P}(X_{t+1}=a_{t+1} | X_{t} = a_t, X_{t-1} = a_{t-1}, \cdots, X_{0} = a_0) =
P(X_{t + 1} = a_{t + 1} | X_t = a_{t})$$
Random walk with absorbing states is not too hard, just solve lots of equations
\subsection{Reflecting Random walk on a line graph}
Suppose we have a random walk $X_t$ takes values on $S = {0, 1, 2, \cdots, N}.$ such that $$P(i, i + 1) = p,
P(i, i - 1) = 1 - p, P(N, N - 1) = P(0, 1) = 1 \text{ and } 0 < p < 1$$
In this case, at the previously absorbing states 0 and 1, they "reflect" so we don't get
stuck and with probability 1 we always return to any state
\subsection{Infinite Line Graph}
Consider the previous example except $S = \mathbb{Z}.$ \par Let $X_t$ describe the
position at time t of the random walk.
\par In this random walk, do we visit any state infinitely many times?
% \begin{answer}
\par \textbf{Answer:} Yes, if and only if $p = 1 - p = \frac{1}{2}$ and the walk is symmetric.
\par \textbf{Extension: } How does this work in $\mathbb{Z}^2$ or $\mathbb{Z}^3$ ??
\par \textbf{Answer to Extension: } $\mathbb{Z}^2$ if symmetric, $\mathbb{Z}^3$ hell
no!(transient random walk)
\subsection{Random walk over a circle}
The transition matrix is 
\[
P=
\begin{pmatrix}
0 & p & 0 & \cdots & 0 & q\\
q & 0 & p & \ddots & \vdots & 0\\
0 & q & 0 & \ddots & 0 & \vdots\\
\vdots & \ddots & \ddots & \ddots & p & 0\\
0 & \cdots & 0 & q & 0 & p\\
p & 0 & \cdots & 0 & q & 0
\end{pmatrix}_{(N+1)\times(N+1)},
\qquad p+q=1.
\]
% \end{answer}
\section{definitions and notation}
\begin{definition}[First definition]
Write the definition and one sentence explaining why it matters.
\end{definition}

\section{Claims, theorems, and proof sketches}
\begin{claim}
State the claim in one precise sentence.
\end{claim}

\begin{proofsketch}
Record the key idea, the first useful equation, and the step that still needs checking.
\end{proofsketch}

\section{Examples and calculations}
\begin{example}[Lecture example]
Write the setup, the calculation, and the conclusion.
\end{example}

\section{Questions and follow-up}
\begin{itemize}
  \item What is still unclear?
  \item Which exercise or reading should I revisit?
  \item TODO: 
\end{itemize}

\end{document}
